\begin{ouverture}[cercle]{16}{Triangle rectangle et cercle circonscrit}{Partie 4 · Géométrie et espace}
{Demi-cercle · Centre · Médiane · Rectangles sans équerre}
\objectif{Caractériser un triangle rectangle par son cercle circonscrit.}
\objectif{Déterminer le centre du cercle circonscrit d’un triangle rectangle.}
\objectif{Construire un angle droit ou un rectangle sans équerre.}
\prerequis{Médiatrices, droite des milieux et théorème de Pythagore (chapitres 12, 14 et 15).}
\end{ouverture}
\pagestyle{otMain}\essentiel\fontsize{11}{15}\selectfont
\phantomsection\label{cours-16-1}
\essSection{1}{Le cercle d’un triangle rectangle}
\begin{theoreme}Si un triangle est rectangle, alors le centre de son cercle circonscrit est le milieu de son hypoténuse : le triangle est inscrit dans un demi-cercle dont l’hypoténuse est un diamètre. Ainsi, la médiane issue du sommet de l’angle droit mesure la moitié de l’hypoténuse.\end{theoreme}
\label{fig-cours-16-1-cours-cercle}\figureGeo{cours-cercle}

\begin{demonstration}ABC est rectangle en A ; I est le milieu de $[BC]$ et J celui de $[AB]$. D’après la droite des milieux, $(IJ)\parallel(AC)$ ; comme $(AC)\perp(AB)$, on a $(IJ)\perp(AB)$. La droite $(IJ)$ est donc la médiatrice de $[AB]$, d’où $IA=IB$. Comme $IB=IC$, I est équidistant des trois sommets : c’est le centre du cercle circonscrit.\end{demonstration}
\begin{exemple}ABC est rectangle en A avec $BC=10$ cm. Le cercle circonscrit a pour centre le milieu I de $[BC]$ et pour rayon 5 cm ; la médiane $[AI]$ mesure 5 cm.\end{exemple}
\clearpage
\phantomsection\label{cours-16-2}
\essSection{2}{Reconnaître un triangle rectangle}
\begin{theoreme}[Réciproque]Si un triangle est inscrit dans un cercle qui a pour diamètre l’un de ses côtés, alors ce triangle est rectangle ; l’angle droit est opposé à ce diamètre. Autrement dit, si la médiane issue d’un sommet mesure la moitié du côté opposé, le triangle est rectangle en ce sommet.\end{theoreme}
\begin{demonstration}A est sur le cercle de diamètre $[BC]$ et de centre O. Le demi-tour de centre O transforme A en A’. Les diagonales $[BC]$ et $[AA']$ de ABA’C ont le même milieu O et la même longueur, celle d’un diamètre : ABA’C est un rectangle, donc $\widehat{BAC}=90^\circ$.\end{demonstration}
\begin{savoirfaire}[Construire sans équerre]Pour un rectangle : tracer un cercle et deux de ses diamètres $[AC]$ et $[BD]$ ; ABCD a des diagonales de même milieu et de même longueur, c’est un rectangle. Pour un angle droit en A sur une droite $(d)$ : un cercle passant par A recoupe $(d)$ en B ; si $[BC]$ est un diamètre, $(AC)\perp(d)$.\end{savoirfaire}
\begin{attention}Un triangle inscrit dans un cercle n’est rectangle que si l’un de ses côtés est un diamètre. Le codage ou l’énoncé doit le garantir ; l’aspect du dessin ne suffit pas.\end{attention}

\clearpage\rubrique{Découverte}
\begin{blocLivre}{ex-16-1}\begin{exercice}[1]{Trouver le centre}
ABC est rectangle en A et $BC=9$ cm. Où se trouve le centre de son cercle circonscrit ? Quel est le rayon de ce cercle ?
\label{fig-ex-16-1}\figureGeo{cer-01}
\end{exercice}\end{blocLivre}
\begin{blocLivre}{ex-16-2}\begin{exercice}[1]{Une médiane}
ABC est rectangle en A, $BC=13$ cm et I est le milieu de $[BC]$. Calculer AI.
\label{fig-ex-16-2}\figureGeo{cer-02}
\end{exercice}\end{blocLivre}
\begin{blocLivre}{ex-16-3}\begin{exercice}[1]{Un angle droit garanti}
Les points B et C sont les extrémités d’un diamètre d’un cercle ; A est un autre point de ce cercle. Que peut-on dire du triangle ABC ?
\label{fig-ex-16-3}\figureGeo{cer-03}
\end{exercice}\end{blocLivre}
\begin{blocLivre}{ex-16-4}\begin{exercice}[1]{Le rayon par Pythagore}
ABC est rectangle en A, avec $AB=6$ cm et $AC=8$ cm. Calculer le rayon de son cercle circonscrit.
\label{fig-ex-16-4}\figureGeo{cer-04}
\end{exercice}\end{blocLivre}
\begin{blocLivre}{ex-16-5}\begin{exercice}[1]{Vrai ou faux ?}
Affirmation : « Tout triangle inscrit dans un cercle est rectangle. » Est-elle vraie ?
\label{fig-ex-16-5}\figureGeo{cer-05}
\end{exercice}\end{blocLivre}

\rubrique{Entraînement}
\begin{blocLivre}{ex-16-6}\begin{exercice}[2]{Une médiane qui décide}
Dans le triangle ABC, I est le milieu de $[BC]$, $BC=7$ cm et $AI=3{,}5$ cm. Démontrer que ABC est rectangle.
\label{fig-ex-16-6}\figureGeo{cer-06}
\end{exercice}\end{blocLivre}
\begin{blocLivre}{ex-16-7}\begin{exercice}[2]{Calculer un angle}
$[AB]$ est un diamètre d’un cercle et C un point de ce cercle, distinct de A et B. On sait que $\widehat{ABC}=32^\circ$. Calculer $\widehat{BAC}$.
\label{fig-ex-16-7}\figureGeo{cer-07}
\end{exercice}\end{blocLivre}
\begin{blocLivre}{ex-16-8}\begin{exercice}[2]{Un rectangle sans équerre}
On trace un cercle de centre O et de rayon 4 cm, puis deux diamètres $[AC]$ et $[BD]$ non perpendiculaires. Démontrer que ABCD est un rectangle et donner la longueur de ses diagonales.
\label{fig-ex-16-8}\figureGeo{cer-08}
\end{exercice}\end{blocLivre}
\begin{blocLivre}{ex-16-9}\begin{exercice}[2]{Des triangles rectangles en M}
A et B sont deux points tels que $AB=6$ cm. Où placer un point M pour que le triangle AMB soit rectangle en M ? Justifier.
\label{fig-ex-16-9}\figureGeo{cer-09}
\end{exercice}\end{blocLivre}
\begin{blocLivre}{ex-16-10}\begin{exercice}[2]{Avec Pythagore}
$[AB]$ est un diamètre d’un cercle de rayon 5 cm, et C un point du cercle tel que $AC=6$ cm. Calculer BC.
\label{fig-ex-16-10}\figureGeo{cer-10}
\end{exercice}\end{blocLivre}

\rubrique{Problèmes}
\begin{blocLivre}{ex-16-11}\begin{exercice}[3]{Une perpendiculaire au compas}
A est un point d’une droite $(d)$. On choisit un point O hors de $(d)$ ; le cercle de centre O passant par A recoupe $(d)$ en B. On trace le diamètre $[BC]$. Démontrer que $(AC)$ est perpendiculaire à $(d)$.
\label{fig-ex-16-11}\figureGeo{cer-11}
\end{exercice}\end{blocLivre}
\begin{blocLivre}{ex-16-12}\begin{exercice}[3]{Dans un demi-cercle}
Un demi-cercle a pour diamètre $[AB]$ avec $AB=10$ cm. Le point C du demi-cercle vérifie $AC=BC$. Calculer AC, exactement puis au centième, et l’aire du triangle ABC.
\label{fig-ex-16-12}\figureGeo{cer-12}
\end{exercice}\end{blocLivre}
\begin{blocLivre}{ex-16-13}\begin{exercice}[3]{Distance du centre à un côté}
ABC est rectangle en A, avec $AB=5$ cm et $AC=12$ cm. O est le centre de son cercle circonscrit et J le milieu de $[AB]$. Calculer le rayon du cercle et la longueur OJ.
\label{fig-ex-16-13}\figureGeo{cer-13}
\end{exercice}\end{blocLivre}
\begin{blocLivre}{ex-16-14}\begin{exercice}[3]{Un rectangle inscrit}
Un rectangle ABCD est inscrit dans un cercle de rayon 6,5 cm, et $AB=5$ cm. Calculer BC puis l’aire du rectangle.
\label{fig-ex-16-14}\figureGeo{cer-14}
\end{exercice}\end{blocLivre}
\begin{blocLivre}{ex-16-15}\begin{exercice}[3]{Sans mesurer d’angle}
Sur un cercle de centre O et de rayon 4 cm, on place B et C tels que O appartienne à $[BC]$, puis un point A du cercle tel que $AB=5$ cm. Calculer AC exactement, puis au centième. Pourquoi n’a-t-on pas besoin de mesurer un angle ?
\label{fig-ex-16-15}\figureGeo{cer-15}
\end{exercice}\end{blocLivre}

\clearpage
